\section{Justificativa de alinhamento}
\label{sec:justificativa}

\subsection{Comunicação}
\begin{enumerate}
\item \canal{WhatsApp e grupos de bairro}: alcança tutores de baixa renda com baixo custo de mídia e alta familiaridade. Apoia a orientação confiável e o boca a boca, mas deve usar mensagens curtas, linguagem simples e consentimento para contato.
\item \canal{Parcerias com ONGs e clínicas}: associa o PetCuida a instituições que já possuem confiança e conhecimento do território. É adequado para apresentar a proposta de créditos solidários e prevenção a tutores que já buscam ajuda.
\item \canal{Instagram/Facebook}: permite explicar a proposta por conteúdo visual, divulgar histórias de impacto e testar públicos por região. A viabilidade é alta porque a produção inicial pode ser feita organicamente pelo grupo e pelos parceiros.
\end{enumerate}

\subsection{Distribuição e entrega}
\begin{enumerate}
\item \canal{Aplicativo/webapp}: entrega o núcleo de valor do PetCuida -- calendário automatizado, prontuário digital, triagem orientativa e busca de parceiros -- em uma experiência integrada.
\item \canal{WhatsApp para onboarding e lembretes}: funciona como porta de entrada para quem não instala um aplicativo ou possui aparelho com pouco espaço. Deve complementar, e não substituir, a avaliação veterinária.
\item \canal{Clínicas veterinárias parceiras}: transforma a descoberta em serviço real, com possibilidade de agenda local, preços sociais e pagamento híbrido. A expansão pode começar com poucos parceiros e uma área geográfica delimitada.
\end{enumerate}

\subsection{Relacionamento}
\begin{enumerate}
\item \canal{Atendimento humano por WhatsApp}: reduz insegurança diante de sintomas e encaminha situações de risco para atendimento profissional. O protocolo deve deixar explícito que a triagem é orientativa e não diagnóstica.
\item \canal{Notificações de prevenção}: materializa o benefício de agir antes da emergência, reforçando vacinação, vermifugação e retornos. A frequência precisa ser controlada para evitar fadiga e descadastros.
\item \canal{Comunidade com ONGs e tutores}: cria escuta contínua, acolhimento e compartilhamento de informação confiável. Também ajuda a identificar barreiras não previstas e a melhorar o produto com evidências do território.
\end{enumerate}
